Most equity, futures and increasingly many fixed-income instruments trade on
\emph{order-driven} markets.
In such a market there is no designated counterparty standing ready to quote a price.
Instead,
the participants themselves supply the prices,
by submitting orders to a venue that stores them and matches them against one another.

An order is an instruction to buy or to sell a given quantity of an instrument.
It carries a \emph{direction} ---
buy, also called \emph{bid}, or sell, also called \emph{ask} ---
and a \emph{size},
that is the number of units the sender is willing to transact.
Orders come in two kinds,
and the distinction between them is the single most useful thing to understand about
the mechanics of trading.

A \emph{limit order} carries, in addition, a price:
it is an instruction to transact at that price or better,
and never worse.
If no counterpart is available at the moment of submission,
the order is not executed.
It is stored by the venue and waits.
A limit order therefore \emph{supplies} liquidity:
it stands in the market as an offer that somebody else may take.

A \emph{market order} carries no price.
It is an instruction to transact immediately,
against whatever the venue currently stores,
at whatever price that turns out to be.
A market order therefore \emph{consumes} liquidity:
it removes offers that somebody else had placed.

The asymmetry between the two is the source of nearly every phenomenon in this chapter.
The sender of a limit order chooses the price but not the time of execution;
the sender of a market order chooses the time but not the price.
Each is exposed to a different risk,
and the trade-off between the two is what an execution algorithm negotiates.

\begin{remark}
\label{remark.tradingIsAnAuction}
An order-driven market is a continuous double auction.
It is continuous because orders may arrive at any instant,
and it is double because both sides of the market bid competitively.
The venue itself takes no position and expresses no view:
it applies a matching rule, and nothing more.
\end{remark}

Prices in such a market are quantised.
A venue defines a \emph{tick size} $\tickSizeOfLOB$,
the smallest permissible price increment,
and every limit order must be priced on the resulting grid $\tickSizeOfLOB \Z$.
The tick size is not a detail:
it sets the granularity at which competition on price can happen,
and when it is large relative to typical price movements it changes
the character of the market entirely.
